% SEGMENT OLP-0054-S01
\documentclass[../../../include/open-logic-section]{subfiles}

\begin{document}

\olfileid{sfr}{infinite}{card-sb}

\olsection{இணைப்பு: ஷ்ரோடர்--பெர்ன்ஸ்டைன் தேற்றத்தை நிறுவுதல்}

முறைசாராத கணக்கோட்பாட்டிலிருந்து விடைபெறும்முன், கடைசியாக ஒரு
முறைசாராத (ஆனால் நுணுக்கமான!{}) நிறுவலைக் கருத வேண்டும்.
இது ஷ்ரோடர்--பெர்ன்ஸ்டைன் தேற்றத்தின்
(\olref[sfr][siz][sb]{thm:schroder-bernstein}) நிறுவல்:
$\cardle{A}{B}$ மற்றும் $\cardle{B}{A}$ எனில் $\cardeq{A}{B}$;
அதாவது, $f \colon A \to B$, $g \colon B \to A$ என்ற
!!{injection}s கொடுக்கப்பட்டால், $h \colon A \to B$ என்ற
!!a{bijection} உள்ளது.

இந்த அத்தியாயத்தில் டெடெகிண்டின் \emph{மூடல்கள்} என்ற கருத்தைப்
பின்பற்றினோம். உண்மையில், இந்தக் கருத்தைப் பயன்படுத்தி டெடெகிண்ட்
ஷ்ரோடர்--பெர்ன்ஸ்டைன் தேற்றத்திற்கு ஓர் அழகிய நிறுவலைத் தந்தார்;
அதனை இங்கு வழங்குவோம். சற்றே வேறுபட்டாலும் சாரத்தில் இதே போன்ற
விளக்கம் வேண்டுமெனில், இந்த நிறுவல் \citet[pp.~157--8]{Potter2004} ஐ
நெருக்கமாகப் பின்பற்றுகிறது. இணையத்தில் சிறிது தேடுவதும்கூட,
$\sqrt{2}$ இன் விகிதமுறாமையைப் போலவே, \emph{பல} சுவையான,
வேறுபட்ட நிறுவல்கள் உள்ள ஒரு தேற்றம் இது என்று உங்களை நம்பவைக்கும்.

% SEGMENT OLP-0054-S02
\olref[sfr][infinite][dedekind]{Closure} இல் உள்ளதுபோன்ற குறியீட்டைப்
பயன்படுத்தி, ஒவ்வொரு கணம் $B$ க்கும் சார்பு~$f$ க்கும்
பின்வருமாறு கொள்க:
\[
\Closureofunder{f}{B} = \bigcap \Setabs{X}{B \subseteq X
\text{ மற்றும் $X$ ஆனது $f$-மூடப்பட்டது}}
\]
இவ்வாறு வரையறுக்கப்படும் $\Closureofunder{f}{B}$ என்பது $B$ ஐக்
கொண்ட மீச்சிறு $f$-மூடிய கணம்; அதாவது:

% SEGMENT OLP-0054-S03
\begin{lem}\ollabel{Closureprops}
	எந்தச் சார்பு $f$ க்கும், எந்த $B$ க்கும்:
	\begin{enumerate}
		\item\ollabel{Closurehaselem} $B \subseteq \Closureofunder{f}{B}$; மேலும்
		\item\ollabel{Closureclosed} $\Closureofunder{f}{B}$ ஆனது $f$-மூடப்பட்டது; மேலும்
		\item\ollabel{Closuresmallest} $X$ ஆனது $f$-மூடப்பட்டும் $B
		\subseteq X$ எனவும் இருந்தால், $\Closureofunder{f}{B} \subseteq X$.
	\end{enumerate}
\end{lem}

% SEGMENT OLP-0054-S04
\begin{proof}
\olref[sfr][infinite][dedekind]{closureproperties} இல் உள்ளதைப்
போலவே நிறுவலாம்.
\end{proof}

ஷ்ரோடர்--பெர்ன்ஸ்டைன் தேற்றத்தை அடைவதற்கு இன்னும் ஒரு உண்மை தேவை:

% SEGMENT OLP-0054-S05
\begin{prop}\ollabel{sbhelper}
$A \subseteq B \subseteq C$ மற்றும் $A \approx C$ எனில்,
$\cardeq{\cardeq{A}{B}}{C}$.
\end{prop}

% SEGMENT OLP-0054-S06
\begin{proof}
$f \colon C \to A$ என்ற !!a{bijection} கொடுக்கப்பட்டுள்ளதாகக் கொள்க.
$F = \Closureofunder{f}{C \setminus B}$ என வைத்து, சார்பகம் $C$ ஆகக்
கொண்ட ஒரு சார்பு $g$ ஐப் பின்வருமாறு வரையறுக்க:
\[
	g(x) =
	\begin{cases}
			f(x) &\text{$x \in F$ எனில்}\\
			x & \text{மற்றபடி}
	\end{cases}
\]
$g$ ஆனது $C \to B$ என்ற !!a{bijection} எனக் காட்டுவோம்.
அதிலிருந்து $\comp{f^{-1}}{g} \colon A \to B$ என்பது
!!a{bijection} எனப் பின்தொடரும்; இது நிறுவலை நிறைவுசெய்யும்.

% SEGMENT OLP-0054-S07
முதலில், $x \in F$ ஆனால் $y\notin F$ எனில் $g(x) \neq g(y)$ என்று
கூறுகிறோம். இதற்கு மாறாக இருப்பதாக முரண்வழியாகக் கொள்க; அப்போது
$y = g(y) = g(x) = f(x)$. $x \in F$ என்பதாலும்,
\olref{Closureprops} இன்படி $F$ ஆனது $f$-மூடப்பட்டது என்பதாலும்,
$y = f(x) \in F$; இது ஒரு முரண்.

% SEGMENT OLP-0054-S08
இப்போது $g(x) = g(y)$ எனக் கொள்க. மேலுள்ளதன்படி,
$x \in F$ எனின், அப்போதும் அப்போது மட்டுமே, $y \in F$.
$x, y \in F$ எனில், $f(x) = g (x) = g(y) = f(y)$;
$f$ ஒரு !!a{bijection} என்பதால் $x = y$.
$x, y \notin F$ எனில், $x = g(x) = g(y) = y$.
ஆகவே $g$ ஓர் !!a{injection}.

% SEGMENT OLP-0054-S09
$\ran{g} = B$ என இன்னும் காட்ட வேண்டும். ஆகவே
$x \in B \subseteq C$ ஐ நிலைநிறுத்துக. $x \notin F$ எனில்,
$g(x) = x$. $x \in F$ எனில், ஏதோ $y \in F$ க்கு $x = f(y)$;
ஏனெனில் அவ்வாறு இல்லாவிட்டால் $F \setminus \{x\}$ ஆனது
$f$-மூடப்பட்டு $C\setminus B$ ஐ உள்ளடக்கியிருக்கும்; இது
\olref{Closureprops} இன்படி இயலாதது. இப்போது $g(y) = f(y) = x$.
\end{proof}

% SEGMENT OLP-0054-S10
இறுதியாக, முதன்மை முடிவின் நிறுவல் இதோ. ஒரு சார்பு $h$ உம் கணம்
$D$ உம் கொடுக்கப்பட்டால், $\funimage{h}{D} = \Setabs{h(x)}{x
\in D}$ என வரையறுக்கிறோம் என்பதை நினைவுகூர்க.

% SEGMENT OLP-0054-S11
\begin{proof}[ஷ்ரோடர்--பெர்ன்ஸ்டைன் தேற்றத்தின் நிறுவல்]
$f \colon A \to B$, $g \colon B \to A$ ஆகியவை !!{injection}s
ஆக இருக்கட்டும். $\funimage{f}{A} \subseteq B$ என்பதால்,
$\funimage{g}{\funimage{f}{A}} \subseteq g[B] \subseteq A$.
மேலும் $g$, $f$ இரண்டும் உட்செலுத்தல்கள் என்பதால்,
$\comp{f}{g} \colon A \to \funimage{g}{\funimage{f}{A}}$ ஆனதும்
ஓர் !!{injection}; உண்மையில், அதன் வீச்சகத்தை நாம் வரையறுத்த
விதத்தினாலேயே $\comp{f}{g}$ ஓர் !!a{bijection}.
ஆகவே $\cardeq{\funimage{g}{\funimage{f}{A}}}{A}$;
எனவே \olref{sbhelper} இன்படி $h \colon A \to \funimage{g}{B}$ என்ற
!!a{bijection} உள்ளது. மேலும் $g^{-1}$ ஆனது
$\funimage{g}{B} \to B$ என்ற !!a{bijection}.
ஆகவே $\comp{h}{g^{-1}} \colon A \to B$ ஓர் !!a{bijection}.
\end{proof}

\end{document}
